\chapter{Аналитическая часть}

\section{Постановка задачи}
Даны матрица $A$ размера $n \times m$ и матрица $B$ размера $m \times q$ с целочисленными элементами. Требуется вычислить их произведение~--- матрицу $C = A \cdot B$ размера $n \times q$, элементы которой определяются формулой~\cite[с.~172]{laaksonen2024}
\begin{equation}
	\label{eq:mult}
	c_{i,j} = \sum_{k=0}^{m-1} a_{i,k} \cdot b_{k,j}, \quad i = \overline{0, n-1}, \quad j = \overline{0, q-1}.
\end{equation}
Произведение определено только в том случае, когда число столбцов матрицы $A$ равно числу строк матрицы $B$.

Для решения задачи рассматриваются следующие алгоритмы:
\begin{itemize}
	\item стандартный алгоритм умножения матриц;
	\item алгоритм Винограда;
	\item оптимизированный согласно варианту алгоритм Винограда в двух реализациях: с заменой деления на 2 побитовым сдвигом и с заменой присваиваний вида $x = x + y$ составным оператором $x \mathrel{+}= y$.
\end{itemize}

\section{Стандартный алгоритм умножения матриц}
Стандартный алгоритм вычисляет каждый элемент $c_{i,j}$ непосредственно по формуле~\eqref{eq:mult} как скалярное произведение $i$-й строки матрицы $A$ на $j$-й столбец матрицы $B$~\cite[с.~172--173]{laaksonen2024}. Для этого используются три вложенных цикла: по строкам $A$, по столбцам $B$ и по индексу суммирования $k$. На вычисление каждого из $n q$ элементов результата затрачивается $m$ умножений и $m$ сложений, поэтому для квадратных матриц размера $N \times N$ временная сложность алгоритма составляет $O(N^3)$~\cite[с.~173]{laaksonen2024}.

\section{Алгоритм Винограда}
Алгоритм Винограда~\cite{winograd1968} основан на том, что каждый элемент результата является скалярным произведением двух векторов, и это произведение можно переписать так, чтобы часть умножений выполнялась заранее, один раз для каждой строки и каждого столбца. Для векторов $V = (v_0, v_1, v_2, v_3)$ и $W = (w_0, w_1, w_2, w_3)$ справедливо
\begin{equation}
	\label{eq:wino-idea}
	V \cdot W = (v_0 + w_1)(v_1 + w_0) + (v_2 + w_3)(v_3 + w_2) - v_0 v_1 - v_2 v_3 - w_0 w_1 - w_2 w_3.
\end{equation}
Слагаемые $v_0 v_1 + v_2 v_3$ зависят только от строки матрицы $A$, а $w_0 w_1 + w_2 w_3$~--- только от столбца матрицы $B$, поэтому их можно вычислить заранее. Для каждой строки $A$ и каждого столбца $B$ вычисляются
\begin{align}
	\label{eq:mulH}
	mulH_i &= \sum_{k=0}^{\floor{m/2}-1} a_{i,2k} \cdot a_{i,2k+1}, \quad i = \overline{0, n-1}, \\
	\label{eq:mulV}
	mulV_j &= \sum_{k=0}^{\floor{m/2}-1} b_{2k,j} \cdot b_{2k+1,j}, \quad j = \overline{0, q-1},
\end{align}
после чего элементы результата находятся по формуле
\begin{equation}
	\label{eq:wino-c}
	c_{i,j} = -mulH_i - mulV_j + \sum_{k=0}^{\floor{m/2}-1} \left(a_{i,2k} + b_{2k+1,j}\right) \cdot \left(a_{i,2k+1} + b_{2k,j}\right).
\end{equation}
При нечётном $m$ последний элемент строки и столбца в пары не попадает, поэтому к каждому элементу результата дополнительно прибавляется слагаемое $a_{i,m-1} \cdot b_{m-1,j}$.

В основном цикле на каждую пару элементов выполняется одно умножение вместо двух в стандартном алгоритме, то есть число умножений в основной части уменьшается примерно вдвое ценой увеличения числа сложений и дополнительной памяти под массивы $mulH$ и $mulV$ размера $n$ и $q$. Асимптотическая сложность алгоритма остаётся равной $O(N^3)$.
